\question[3]
Gegeben ist der folgende Python-Code:

\begin{lstlisting}
x = 1

def f():
    x = 2
    print(x)

f()
print(x)
\end{lstlisting}

\begin{parts}
\part Gib an, was das Programm ausgibt.
\part Erläutere die Ausgabe.
\end{parts}

\begin{solutionbox}{4.5cm}
\begin{parts}
\part \texttt{2} und danach \texttt{1}
\part Durch die Wertzuweisung \texttt{x = 2} in der Funktion entsteht eine neue lokale Variable \texttt{x}.
Sie hat zufällig denselben Bezeichner wie die globale Variable, ist aber eine andere Variable.
In der Funktion wird die lokale Variable ausgegeben, außerhalb die globale. Die globale Variable bleibt unverändert.
\end{parts}
\end{solutionbox}
